\section{Introduction}

This paper records the public version of a closed audit chain: finite tail
graphs, cycle realizability, phase-decomposed renewal drift, and
Foster-Lyapunov residue diagnostics.  It is in the genre of computational
mathematics and experimental mathematics.  It makes no claim of proving the
Collatz conjecture, no descent theorem, and no theorem about deterministic
per-orbit descent.

The object is the accelerated odd Collatz map
\[
  \Syr(n)=\frac{3n+1}{2^{\vTwo(3n+1)}} ,
\]
applied to odd positive integers.  The project began as a certificate-search
framework and ended as a reproducible finite diagnostic framework.  The audit
history in \doclink{PROJECT_JOURNEY.md}{PROJECT\_JOURNEY.md} records several
demoted claims: a per-step running-debt Lyapunov ansatz failed at roughly
50 percent non-decrease, a structural Diophantine exclusion was a resolution
artifact, and an early closed-form reading of the renewal constant was rejected.

The current finite picture is narrower.  On tested LTE-closed tail graphs,
bounded integer-realizable simple cycles are filtered down to the elementary
\([2]\) cycle, whose factor is \(3/4\).  Independently, the phase-decomposed
renewal drift audit finds the same value as a mean-step target after separating
tail-internal and post-exit phases.  Finally, the Foster audit finds that
\[
  V(n)=\log_2(n)+\vTwo(n+1)
\]
does not give one-step residue-uniform negativity, but does give an \(m\)-step
finite residue diagnostic at \(m=16\) on the tested windows.

There are now two relevant Chang 2026 papers.  The human-LLM collaboration
paper \cite{chang2026collatz} is a methodological and structural companion,
while the later structural reduction \cite{chang2026reduction} reduces the
Collatz question to a fixed-modulus, one-bit orbit-mixing problem.  The
compatibility analysis in
\artifact{chang_2603_25753_compatibility.md}{chang\_2603\_25753\_compatibility.md}
reads the two efforts as addressing the same distributional-to-pointwise wall
from complementary angles: Chang isolates a one-bit balance bottleneck, and
this framework supplies an \(m=16\) Foster residue diagnostic that can be
composed with that reduction as a joint argument.

\begin{caveat}
All quantitative statements below are tied to saved JSON artifacts.  The
language ``finite empirical diagnostic'' is intentional: the computations are
finite, sampled or bounded, and do not imply a global statement for all
positive integers.
\end{caveat}
